A design is an incidence structure of points and blocks. 
The incidence matrix of a design has rows corresponding to the points 
and columns corresponding to the blocks. 
An entry in a certain row and column is one if and only if the point associated with the row 
is contaned in the block associated with the column, zero otherwise.
A decomposition of the design is a partition of the points and blocks such that 
each class consists either exclusively of points or exclusively of blocks. 

\bigskip


A decomposition is point-tactical if for all points, 
the number of incident lines in the $j$th block class depends only on the class of the point.
It the point belongs to class $i,$ this number is denoted as $a_{ij}.$
A decomposition is block-tactical if for all blocks, 
the number of incident points in the $i$th point class depends only on the class of the block.
If the block belongs to class $j,$ this number is denoted as $b_{ij}.$

\bigskip

A projective plane of order $n$ is a 
design with $n^2+n+1$ points and equally 
many blocks (also called lines), each of size $n+1$ such that 
any two points lie in exactly one block and any two blocks have exactly one point in common.
Projective planes are known to exist 
for all $n = q$ which are a power of a prime. 
This follows from a construction which utilizes the projective geometry $\PG(2,q).$ 
Points are the one-dimensional subspaces of $\bbF_q^3,$
blocks are the two-dimensional subspaces of $\bbF_q^3,$ and
incidence is natural (inclusion of subspaces). 
The automorphism group of this design is the collineation group of the projective space.
Projective planes other than these exist, though none are known when $n$ is not a prime power.
The number of lines through a point equals the number of points on a line. 
The fact that these numbers exist imply 
that there is a tactical decomposition. 
Namely, the trivial decomposition with two classes, 
one containing all points and one containing all lines.
The structure constants of the decomposition are the numbers just described.

\bigskip

Table~\ref{tab:design:create} 
lists Orbiter commands to create designs.
\orbitertableofcommands{design_create}
Table~\ref{tab:design:activity}
lists Orbiter activities for designs.
\orbitertableofcommands{design_activity}

The command 

\sourcecode{design_PG_2_3}

creates the design $\PG(2,3).$ 
\orbiteroutput{GRAPHICS/LATEX/design_PG_2_3}
The blocks of the design are encoded in the lexicographic ordering of $k$-subsets (here $k=4$).
The program also displays the tactical decomposition scheme of the design, which is
\orbiteroutput{GRAPHICS/LATEX/design_PG_2_3_td}
In Section~\ref{sec:canonical:form:objects:design:theory}, we will show how to compute 
further properties of the design.


\bigskip

The command 

\sourcecode{wreath_product_designs_n4_k2_inc.txt}

creates a design on $8$ points invariant under the wreath product $\Sym(4) \wr \Sym(2).$
The design has 12 blocks of size 4. 
The command 

\sourcecode{wreath_product_designs_n8_k6_inc.txt}

creates a design on $16$ points invariant under 
the wreath product $\Sym(8) \wr \Sym(2).$
The design has 3920 blocks of size 6. 
We will compute the automorphism groups of these 
two designs in Section~\ref{sec:canonical:form:objects:incidence:geometries}.



\bigskip

A design can be created by listing the blocks in coded fashion. 
The codes correspond to the 
ranking of $k$-subsets of the $v$ set which is the underlying set of the design. 
The following command creates a configuration $15_4 20_3$ from its set of blocks. 

\sourcecode{design_20}




\bigskip

It is also possible to create a design by listing the blocks as sets. 
Here is an example. At first, we create a makefile variable 
for a csv file containing the blocks:

\sourcecodevariable{GEO_BLOCKS_600}

The following command creates the csv file from the makefile variable using the unix echo command. 
The sed command is necessary to remove space characters that are inserted by the echo command.
The Orbiter command then reads the list of blocks from the csv file. 
After that, it creates the design from the list of blocks. 
As we will see below, the design has an automorphism group of order 600 
and is related to the Latin Square of order 5 associated with the cyclic group of order 5. 

\sourcecode{design_600}

The command also computes the intersection matrix of the design. 
The design is a transversal design. 
The configuration graph is three copies of the complete graph $K_5.$


\bigskip


How can we convince ourselves that the design 
created previously is indeed associated to a certain Latin Square?
There are exactly two Latin Squares of order 5. 
One is the Cayley table of the cyclic group of order 5, 
the other is the multiplication table of the 
non-associative loop of order 5, see Table~\ref{tab:lsq5}. 

\input GRAPHICS/TABLES/OTHER/lsq5.tex

We create makefile variables for the two tables:

\sourcecodevariable{LSQ_5A_TABLE}

\sourcecodevariable{LSQ_5B_TABLE}

The next two commands create the associated linear spaces:

\sourcecode{LSQ_5A}

\sourcecode{LSQ_5B}

Finally, we perform the canonical form algorithm on these two linear spaces. 
We note that any isomorphism between 
the Latin Squares corresponds to an isomorphism 
between the associated linear spaces.  
Likewise, any automorphism of a given Latin Square 
corresponds to an automorphism of the 
associated linear space.

\sourcecode{LSQ_5_c}

This establishes the fact that the geometries are not isomorphic, 
and it also produces the automorphism groups of each geometry. 
It turns out that the Latin Square of order 5 arising 
from the cyclic group has a symmetry group of order 600. 
The second Latin Square arising from the loop has a symmetry group of order 72. 
Orbiter produces a report listing generators for each group.
For more information about the canonical form of designs, 
see Section~\ref{sec:canonical:form:objects:design:theory}.
Since there is only one linear space of Latin Square 5 type with a group of order 600, 
the object constructed earlier must be that same space 
which is associated to the Cayley table of the cyclic group of order 5.


